\documentclass[12pt,a4paper]{article}

\usepackage{xetexko}
\usepackage{geometry}
\usepackage{xcolor}
\usepackage{fancybox}
\usepackage{array}
\usepackage{booktabs}
\usepackage{tabularx}
\usepackage{fancyhdr}
\usepackage{colortbl}

\geometry{
  a4paper,
  top=20mm, bottom=20mm,
  left=25mm, right=25mm
}

\setmainfont{Apple SD Gothic Neo}
\setsansfont{Apple SD Gothic Neo}
\setmainhangulfont{Apple SD Gothic Neo}[BoldFont=Apple SD Gothic Neo Bold]

\definecolor{headerblue}{RGB}{0, 32, 96}
\definecolor{sectiongray}{RGB}{217, 217, 217}
\definecolor{lightgray}{RGB}{242, 242, 242}

\setlength{\parindent}{0pt}
\setlength{\parskip}{3pt}

\pagestyle{fancy}
\fancyhf{}
\fancyfoot[C]{\small\thepage}
\renewcommand{\headrulewidth}{0pt}

% 섹션 헤더 커맨드
\newcommand{\sectionheader}[2]{%
  \vspace{3mm}%
  \noindent\colorbox{sectiongray}{%
    \parbox{\dimexpr\textwidth-2\fboxsep\relax}{%
      \vspace{2pt}%
      \textbf{#1}\\%
      {\small #2}%
      \vspace{2pt}%
    }%
  }%
  \vspace{2mm}%
}

\newcommand{\sectionheadersimple}[1]{%
  \vspace{3mm}%
  \noindent\colorbox{sectiongray}{%
    \parbox{\dimexpr\textwidth-2\fboxsep\relax}{%
      \vspace{3pt}%
      \textbf{#1}%
      \vspace{3pt}%
    }%
  }%
  \vspace{2mm}%
}

\begin{document}

% ── 제목 박스 ────────────────────────────────────────────
\noindent\colorbox{headerblue}{%
  \parbox{\dimexpr\textwidth-2\fboxsep\relax}{%
    \vspace{6pt}%
    \centering{\large\bfseries\color{white}
    『2026 관광데이터 활용 공모전』- ① 웹·앱 개발 부문 제안서}%
    \vspace{6pt}%
  }%
}

\vspace{2mm}

% ── 작성 안내 박스 ───────────────────────────────────────
\noindent\fbox{%
  \parbox{\dimexpr\textwidth-2\fboxsep-2\fboxrule\relax}{%
    \vspace{3pt}%
    {\small
    - 작성항목 및 양식 수정 금지, \textbf{모든 항목 작성 필수}, 이미지(사진), 표 등 삽입 가능\\
    - 분량: \textbf{5페이지 이내}로 작성(이미지, 표 등 일체 포함), 글자 포인트는 \textbf{12포인트 이상}\\
    - 제출파일 형식: \textbf{PDF(10MB 미만)}%
    }%
    \vspace{3pt}%
  }%
}

\vspace{3mm}

% ══════════════════════════════════════════════════════════
% 1) 서비스 기획배경 및 필요성
% ══════════════════════════════════════════════════════════
\sectionheader{1) 서비스 기획배경 및 필요성}{※ 서비스를 기획·구상하고 제안한 배경과 필요성 작성}

\textbf{1. 서비스 기획 배경}

제주도는 연간 1,500만 명 이상이 방문하는 국내 최대 관광지이지만, 현행 관광 앱의 대부분은 맛집·숙박·교통 중심의 정보 제공에 머물고 있다. 제주특별자치도는 설화 182건·민담 323건 등 총 505건의 무형문화유산 콘텐츠를 공공 데이터로 개방하고 있음에도, 이를 실제 관광 경험과 연결한 서비스는 존재하지 않는다.

탐라담은 이 공백을 채운다. 사용자가 설화 테마·여행 지역·여행 일자를 선택하면 AI가 \textbf{설화 밀도를 기준으로 코스를 선별}하여, 이야기와 연결된 장소들로 여행이 구성된다. 기존 관광 앱이 명소를 먼저 고르고 정보를 붙이는 방식이라면, 탐라담은 설화가 먼저 특정되고 그 장소로 여행자를 데려간다. LLM 및 RAG 기술의 발전으로 이를 실용 수준에서 구현할 수 있는 시점에 도달했다고 판단하였다.

\vspace{2mm}
\textbf{2. 서비스 필요성}

\begin{itemize}
  \item \textbf{무형문화유산 접근성 제고}: 공공 데이터로 개방된 505건의 설화·민담이 관광 서비스와 연결되지 않아 일반인에게 알려지지 않고 있다. 여행이라는 직접 경험과 결합함으로써 문화유산의 실질적 활용도를 높인다.
  \item \textbf{차별화된 관광 경험 제공}: 정보 소비형 여행에서 이야기 속 주인공이 되는 몰입형 여행으로 패러다임을 전환하여, 재방문 동기와 체류 시간을 늘리는 관광 콘텐츠를 제공한다.
  \item \textbf{공공 데이터 연계 가치 창출}: 한국관광공사 TourAPI, 제주관광공사 Visit Jeju 코스 데이터, 제주특별자치도 설화·민담 API를 AI로 결합하여, 개별 데이터로는 불가능한 새로운 관광 서비스를 실현한다.
\end{itemize}

\vspace{2mm}

% ══════════════════════════════════════════════════════════
% 2) 서비스 개요
% ══════════════════════════════════════════════════════════
\sectionheader{2) 서비스 개요}{%
  ※ 기획 서비스에 대한 소개 및 구현 예정인 주요 기능 위주로 작성\newline
  ※ 지역 특화 서비스의 경우, 특화 지역 관련 내용 작성 필수}

\textbf{1. 기획 서비스 소개}

제주 설화·민담 505건을 AI로 분석해, 이야기 속 배경 장소로 여행 코스를 역설계하는 iOS 여행 앱 \textbf{「탐라담(耽羅談)」}

\vspace{2mm}
\textbf{2. 기획 서비스 주요 기능}

\begin{itemize}
  \item \textbf{설화 테마 기반 AI 코스 추천}: '신비로운 제주', '바다의 제주' 등 설화 테마를 선택하면 LangGraph AI 에이전트가 Visit Jeju 9,134개 실제 관광객 코스 중 설화 밀도가 높은 코스를 선별해 추천한다.
  \item \textbf{GPS 탐험 모드}: iOS CoreLocation 지오펜스로 장소 도착을 자동 감지한다. 도보 반경 100m·차량 반경 300m 진입 후 30초 체류 시 도착으로 판정하며, AI 동행자가 자동으로 말을 건다.
  \item \textbf{AI 설화 동행자 채팅}: 마을 할망·돌하르방·도깨비·해녀 중 원하는 동행자를 선택하면, RAG 기반으로 해당 장소 주변 설화를 실시간 검색해 SSE 스트리밍으로 이야기를 들려준다.
  \item \textbf{여행 일지 자동 생성 및 공유}: 탐험 종료 후 GPT-4o가 방문 이력·동행자 대화 맥락을 바탕으로 여행자 1인칭 산문 일지를 자동 생성하며, iOS ShareLink로 SNS 공유가 가능하다.
  \item \textbf{앱 세션 복원}: 탐험 중 앱이 강제 종료되어도 재시작 시 진행 중인 여행 세션을 자동 복원하여 여행의 연속성을 보장한다.
\end{itemize}

\vspace{2mm}
\textbf{3. 서비스 차별성}

\begin{center}
\begin{tabular}{|>{\bfseries}p{3.2cm}|p{4.8cm}|p{5.8cm}|}
\hline
\rowcolor{lightgray}
\textbf{항목} & \textbf{기존 관광 앱} & \textbf{탐라담} \\
\hline
콘텐츠 방향 & 명소·맛집·숙박 중심 (설화 연결 없음) & \textbf{설화 밀도 기준 코스 선별 → 이야기 속 장소로 안내} \\
\hline
여행 경험 & 정보 조회 중심 & \textbf{이야기 속 주인공이 되는 몰입형 경험} \\
\hline
AI 동행자 & 없음 & \textbf{RAG 기반 설화 해설, 제주 방언 대화} \\
\hline
데이터 결합 & 단일 출처 & \textbf{공공 API 3종 AI 교차 결합} \\
\hline
\end{tabular}
\end{center}

\vspace{2mm}
\textbf{4. 서비스 내 지역 특화 관련 사항 (제주 특화)}

\begin{itemize}
  \item \textbf{제주 전용 설화 데이터}: 제주특별자치도 공공 API(E07 설화정보·E08 민담정보)로 수집한 설화 182건·민담 323건 전량을 ChromaDB에 벡터화하여 서비스 핵심 콘텐츠로 활용한다.
  \item \textbf{제주 방언 AI 동행자}: 동행자가 "혼저 옵서, 삼춘", "하영 반갑수다" 등 제주 토박이 표현을 자연스럽게 구사하며 현지 몰입감을 제공한다.
  \item \textbf{제주 고유 문화 캐릭터}: 마을 할망(마을 수호신), 돌하르방(제주 수호석상), 도깨비, 해녀 등 제주 문화에 기반한 4종 페르소나를 제공하여 제주다운 여행 경험을 구현한다.
\end{itemize}

\vspace{2mm}

% ══════════════════════════════════════════════════════════
% 3) 데이터 활용 방안
% ══════════════════════════════════════════════════════════
\sectionheadersimple{3) 데이터 활용 방안 \quad {\normalfont\small *한국관광공사 OpenAPI 활용 필수}}

\textbf{1. 활용 예정 한국관광공사 OpenAPI}

한국관광공사 \textbf{국문 관광정보 서비스 (KorService2)} — 아래 4개 API를 실제 서비스에 활용 중

\begin{center}
\begin{tabular}{|p{4.8cm}|p{8.5cm}|}
\hline
\rowcolor{lightgray}
\textbf{API 명칭} & \textbf{활용 목적} \\
\hline
locationBasedList2 & GPS 반경 500m 기준 관광지 검색 — 장소명·contentId 확보 \\
\hline
detailCommon2 & contentId로 관광지 개요·주소·전화번호 조회 \\
\hline
detailImage2 & contentId로 관광지 사진 최대 5장 조회 \\
\hline
detailIntro2 & 운영시간·휴무일·입장료·주차 정보 조회 \\
\hline
\end{tabular}
\end{center}

\vspace{2mm}
\textbf{2. 데이터 활용 방식}

GPS 탐험 모드에서 사용자가 코스 장소에 도착하면, 장소명·GPS 좌표를 \texttt{locationBasedList2}로 KTO DB에서 매칭하여 contentId를 확보한다. 이어서 \texttt{detailCommon2}·\texttt{detailImage2}·\texttt{detailIntro2}를 호출해 사진·설명·운영정보를 한 화면에 표시한다.

조회 결과는 7일 SQLite 캐시로 관리하여 동일 장소 재방문 시 API 호출 없이 즉시 표시하고, KTO DB에 등록되지 않은 장소는 빈 필드 처리로 앱 안정성을 유지한다. 또한 관광지 개요(overview) 필드를 AI 동행자 채팅 컨텍스트로 주입하여, 설화와 장소 정보가 결합된 풍부한 해설을 생성한다.

\vspace{2mm}

% ══════════════════════════════════════════════════════════
% 4) 서비스 발전 방향
% ══════════════════════════════════════════════════════════
\sectionheader{4) 서비스 발전 방향}{※ 기획 서비스 발전 가능성, 기대효과 등에 대해 작성}

\textbf{1. 개발 서비스 향후 발전방향}

\begin{itemize}
  \item \textbf{Android 플랫폼 확장}: 현재 iOS 전용으로 운영하며, 서비스 트래픽과 사용자 기반이 충분히 확보되는 시점에 Android 버전을 추가 출시할 계획이다.
  \item \textbf{다국어 지원}: 영어·일본어·중국어 동행자 페르소나를 추가하여 외국인 관광객으로 대상을 확장하고, 제주 설화를 글로벌 관광 콘텐츠로 발전시킨다.
  \item \textbf{동행자 페르소나 고도화}: 현재 4종 페르소나를 사용자 피드백 기반으로 지속 확장하고, 각 캐릭터별 고유 어투·설화 해석 스타일을 강화한다.
  \item \textbf{전국 설화 지역 확장}: 제주 전용 서비스에서 출발하여 전국 무형문화유산 공공 데이터로 콘텐츠를 확장, 다른 지역 고유 설화 기반 여행 코스를 제공하는 전국 단위 플랫폼으로 성장을 도모한다.
  \item \textbf{개인화 추천 고도화}: 사용자 방문 이력 및 동행자 대화 패턴을 분석하여 설화 취향을 학습하고, 재방문 시 더 정교한 코스를 자동 추천하는 개인화 기능을 개발한다.
\end{itemize}

\end{document}
